\documentclass[11pt]{article}

\usepackage[utf8]{inputenc}
\usepackage[T1]{fontenc}
\usepackage{geometry}
\usepackage{amsmath}
\usepackage{amssymb}
\usepackage{booktabs}
\usepackage{hyperref}
\usepackage{xurl}
\usepackage{xcolor}
\usepackage{listings}
\usepackage{enumitem}
\usepackage{parskip}
\usepackage{graphicx}
\graphicspath{{../images/}}

\geometry{margin=1in}

\lstset{
  language=Java,
  basicstyle=\ttfamily\small,
  keywordstyle=\color{blue},
  commentstyle=\color{gray},
  stringstyle=\color{teal},
  showstringspaces=false,
  columns=fullflexible,
  frame=single,
  framerule=0.2pt,
  breaklines=true
}

\setlist[itemize]{topsep=0.3em,itemsep=0.2em}
\setlist[enumerate]{topsep=0.3em,itemsep=0.2em}

% Simple per-section source line.
\newcommand{\notesources}[1]{\par\noindent{\small\color{gray}\textit{Sources:} \sloppy #1\par}}

% Unit-wise source strings for this subject (used by slides/notes).
% Path is relative to the lecture `latex/` folder (the compilation CWD).
% OOPS with Java (CSEG1044) — unit-wise sources
%
% These are short, student-facing reference strings intended to be used with:
%   - \slidesources{...}  (in beamer slides)
%   - \notesources{...}   (in lecture notes)
%
% Style inspiration: other subjects in My-Subjects/ use semicolon-separated sources
% and include an "accessed" date for web documentation.

\newcommand{\OOPSUnitISources}{Schildt, \textit{Java: A Beginner's Guide} (10e), McGraw-Hill (Java fundamentals + OOP basics).; Oracle Java Tutorials: Getting Started + Language Basics + Classes and Objects (accessed \today).; Java SE API docs (e.g., \texttt{String}, \texttt{Scanner}) (accessed \today).}

\newcommand{\OOPSUnitIISources}{Schildt, \textit{Java: The Complete Reference} (12e), McGraw-Hill (classes, inheritance, packages, interfaces).; Bloch, \textit{Effective Java} (3e), Addison-Wesley, 2018 (design choices; interfaces vs abstract classes).; Oracle Java Tutorials: Classes and Objects + Interfaces + Packages (accessed \today).; Java SE API docs (e.g., \texttt{Comparable}, \texttt{Runnable}) (accessed \today).}

\newcommand{\OOPSUnitIIISources}{Schildt, \textit{Java: The Complete Reference} (12e) (nested/inner classes, exceptions, threads, I/O).; Oracle Java Tutorials: Nested Classes + Exceptions + Concurrency + Basic I/O (accessed \today).; Java SE API docs (e.g., \texttt{Thread}, \texttt{AutoCloseable}) (accessed \today).}

\newcommand{\OOPSUnitIVSources}{Schildt, \textit{Java: The Complete Reference} (12e) (generics, lambdas, Swing, JDBC).; Bloch, \textit{Effective Java} (3e) (generics + lambdas).; Oracle Java Tutorials: Generics + Lambda Expressions + Swing + JDBC Basics (accessed \today).; Java SE API docs (e.g., \texttt{java.util.function}, \texttt{javax.swing}, \texttt{java.sql}) (accessed \today).}

\newcommand{\OOPSUnitVSources}{Schildt, \textit{Java: The Complete Reference} (12e) (Collections Framework + wrapper classes).; Oracle Java docs: Collections Framework overview (accessed \today).; Java SE API docs (\texttt{java.util} collections; wrapper classes in \texttt{java.lang}) (accessed \today).}

\newcommand{\OOPSUnitVISources}{Git SCM: \textit{Pro Git} (2e) (version control workflow), accessed \today.; GitHub Docs (repositories, issues, pull requests), accessed \today.; Oracle Java Tutorials: Swing + JDBC (accessed \today).; JUnit 5 User Guide (accessed \today).; Apache Maven Project: Getting Started (accessed \today).; Gradle User Manual: Getting Started (accessed \today).; UML class diagrams: UML 2.x overview/spec (accessed \today).}

\newcommand{\OOPSMiscWhyInterfacesSources}{Schildt, \textit{Java: The Complete Reference} (12e) (interfaces).; Bloch, \textit{Effective Java} (3e) (prefer interfaces where appropriate).; Oracle Java Tutorials: Interfaces (accessed \today).; Java SE API docs (e.g., \texttt{Comparable}, \texttt{Runnable}) (accessed \today).}



\title{Object Oriented Programming (CSEG1044)\\Unit 04 -- Lecture 02 Notes\\Generic Methods + Bounds + Wildcards}
\author{Tofik Ali}
\date{\today}

\begin{document}
\maketitle
\tableofcontents

\section{Lecture Overview}
In the previous lecture you learned generic classes like \texttt{Box<T>}.
Now we learn three ideas that appear frequently in real Java code:
\begin{itemize}
  \item \textbf{generic methods}: \texttt{<T>} methods that work for many types,
  \item \textbf{bounds}: restricting which types are allowed (\texttt{extends}),
  \item \textbf{wildcards}: flexible method parameters for collections (\texttt{? extends}, \texttt{? super}).
\end{itemize}
\notesources{\OOPSUnitIVSources}

\section{Syllabus Alignment (Unit 04)}
\textbf{Unit 04:} generic methods; bounded types; wildcards

\section{Learning Outcomes}
After this lecture, you should be able to:
\begin{itemize}
  \item Write a generic method and explain why it reduces duplication.
  \item Apply bounded type parameters (intro level) to restrict allowed types.
  \item Use wildcards with collections for common patterns in real code.
\end{itemize}

\section{Key Terms (Quick)}
\begin{itemize}
  \item \textbf{Generic method}: method with its own type parameter, e.g., \texttt{static <T> T first(List<T> list)}.
  \item \textbf{Bound}: restriction on type parameter, e.g., \texttt{<T extends Number>}.
  \item \textbf{Wildcard}: unknown type parameter, e.g., \texttt{List<?>}.
  \item \textbf{Upper-bounded wildcard}: \texttt{? extends Type} (read/produce).
  \item \textbf{Lower-bounded wildcard}: \texttt{? super Type} (write/consume).
  \item \textbf{Invariance}: \texttt{List<Integer>} is not a subtype of \texttt{List<Number>}.
\end{itemize}

\section{Detailed Notes}
\subsection{Generic methods}
Sometimes only one method needs to be generic (not the whole class).
Example: a method that prints any array.

\textbf{Syntax:}
\begin{itemize}
  \item type parameter list comes before return type: \texttt{public static <T> void print(T[] a)}.
\end{itemize}

\subsection{Bounds (restrict allowed types)}
Without bounds, \texttt{T} can be any reference type.
With bounds, you restrict it:
\begin{itemize}
  \item \texttt{<T extends Number>} means \texttt{T} must be \texttt{Number} or a subclass (\texttt{Integer}, \texttt{Double}, ...).
  \item Common pattern: \texttt{<T extends Comparable<T>>} to support comparisons.
\end{itemize}

\subsection{Wildcards: making collection parameters flexible}
The key beginner confusion is:
\begin{itemize}
  \item \texttt{List<Integer>} is \textbf{not} a \texttt{List<Number>} (invariance).
\end{itemize}
So how do we write a method that can accept \texttt{List<Integer>} \emph{or} \texttt{List<Double>}?
We use wildcards.

\subsubsection{\texttt{? extends T} (producer/read)}
If a method only needs to \textbf{read} from a list, use \texttt{? extends}.
Example: sum of numbers:
\begin{itemize}
  \item \texttt{List<? extends Number>} accepts \texttt{List<Integer>}, \texttt{List<Double>}, ...
\end{itemize}
\textbf{Limitation:} you cannot safely add new elements (except \texttt{null}) because the exact subtype is unknown.

\subsubsection{\texttt{? super T} (consumer/write)}
If a method needs to \textbf{add} elements of type \texttt{T} to a list, use \texttt{? super T}.
Example: add integers to a list that can store integers or any of their supertypes.

\paragraph{PECS (memory trick).} Producer Extends, Consumer Super.
\notesources{\OOPSUnitIVSources}

\section{Worked Example: generic methods + bounds + wildcards}
\subsection*{Generic method (print array)}
\begin{lstlisting}
public class GenericDemo {
    public static <T> void printArray(T[] a) {
        for (T x : a) System.out.print(x + " ");
        System.out.println();
    }
}
\end{lstlisting}

\subsection*{Bounded type parameter (max)}
\begin{lstlisting}
public static <T extends Comparable<T>> T max(T a, T b) {
    return (a.compareTo(b) >= 0) ? a : b;
}
\end{lstlisting}

\subsection*{Upper-bounded wildcard (sum)}
\begin{lstlisting}
import java.util.List;

public static double sum(List<? extends Number> list) {
    double s = 0.0;
    for (Number n : list) s += n.doubleValue();
    return s;
}
\end{lstlisting}

\subsection*{Lower-bounded wildcard (add)}
\begin{lstlisting}
import java.util.List;

public static void addOnes(List<? super Integer> list, int count) {
    for (int i = 0; i < count; i++) list.add(1);
}
\end{lstlisting}

\section{In-Class Exercises (with brief solutions)}
\begin{enumerate}
  \item Why is \texttt{List<Integer>} not a subtype of \texttt{List<Number>}?\par
  \textbf{Answer:} generics are invariant in Java; allowing it would break type safety (you could add a \texttt{Double} into a \texttt{List<Integer>} through a \texttt{List<Number>} reference).

  \item If you only need to read from a list, which wildcard is common: \texttt{extends} or \texttt{super}?\par
  \textbf{Answer:} \texttt{? extends ...}.

  \item What is the PECS rule?\par
  \textbf{Answer:} Producer Extends, Consumer Super.
\end{enumerate}

\section{Self-Study Practice (no solutions)}
\begin{enumerate}
  \item Write a generic method \texttt{swap(T[] a, int i, int j)}.
  \item Write \texttt{min} using \texttt{<T extends Comparable<T>>}.
  \item Write a method that copies from one list to another using wildcards (hint: one produces, one consumes).
\end{enumerate}

\section{Checkpoint Questions (with sample answers)}
\textbf{Q1.} When do you use \texttt{? extends T}?\par
\textbf{Sample answer:} Use \texttt{? extends T} when your method only needs to read values of type \texttt{T} from a structure and doesn't need to add elements, because it can accept multiple subtype lists safely.

\vspace{0.6em}
\textbf{Q2.} What does a bound like \texttt{<T extends Number>} mean?\par
\textbf{Sample answer:} It means the type parameter \texttt{T} must be \texttt{Number} or a subclass of \texttt{Number}, so you can safely use \texttt{Number}-level operations inside the method/class.

\end{document}
